\section{Retroproyección filtrada}

La fórmula de inversión anterior puede escribirse como un procedimiento operativo. Para cada ángulo $\theta$, se filtra la proyección $t\mapsto(\mathcal Rf)(t,\theta)$ en la variable del detector y luego se aplica la retroproyección.\\

El símbolo del filtro rampa es
\[
\tau_{\mathrm{rampa}}(\sigma)=|\sigma|.
\]
Con la notación de multiplicadores del capítulo anterior,
\[
T_{\tau_{\mathrm{rampa}}}h
=
\mathcal F_t^{-1}
\left[
|\sigma|\,\mathcal F_t h
\right].
\]
Para $f\in\mathcal S(\mathbb R^2)$, las proyecciones tienen regularidad suficiente para aplicar este multiplicador en la variable $t$.
Aplicado a cada proyección, esto produce el sinograma filtrado
\[
g_{\mathrm{filt}}(t,\theta)
=
\mathcal F_t^{-1}
\left[
|\sigma|\,\mathcal F_t(\mathcal Rf)(\sigma,\theta)
\right](t).
\]
La fórmula de inversión se convierte entonces en
\[
f
=
\mathcal R^*g_{\mathrm{filt}}
=
\mathcal R^*
\mathcal F_t^{-1}
\left[
|\sigma|\,\mathcal F_t(\mathcal Rf)
\right].
\]
Esta es la forma canónica de la retroproyección filtrada en esta tesis.
Equivalente y abreviadamente,
\[
f=\mathcal R^*T_{\tau_{\mathrm{rampa}}}\mathcal Rf,
\qquad
\tau_{\mathrm{rampa}}(\sigma)=|\sigma|.
\]

El factor $|\sigma|$ incorpora el jacobiano radial del cambio de variables, pero su crecimiento sin cota amplifica componentes de alta frecuencia. Esta observación conecta con la discusión de multiplicadores no acotados del capítulo anterior y motiva el estudio posterior de filtros modificados con símbolos $\tau(\sigma)$.

Escribiendo $g=\mathcal Rf$, la Figura~\ref{fig:fbp-flujo} presenta el procedimiento como una cadena de operaciones. Para un símbolo general $\tau$, la salida se denota $f_\tau$; para la rampa canónica $\tau(\sigma)=|\sigma|$, en el modelo continuo ideal se recupera $f$. En el capítulo computacional, la multiplicación por el símbolo se representará mediante el vector $M_\tau$, o la matriz diagonal asociada, aplicado a cada proyección en frecuencia.

\begin{figure}[htbp]
\centering
\resizebox{0.92\textwidth}{!}{\shorthandoff{<>}
\begin{tikzpicture}[node distance=1.35cm, font=\small]
    \tikzstyle{block}=[draw=black, fill=gray!10, rounded corners=2pt, minimum width=2.25cm, minimum height=0.78cm, align=center]
    \tikzstyle{data}=[align=center, minimum width=2.1cm]
    \node[data] (sino) {sinograma\\$g=\mathcal Rf$};
    \node[block, right=of sino] (fft) {Fourier en $t$\\$\mathcal F_t$};
    \node[data, right=of fft] (ghat) {$\widehat g(\sigma,\theta)$};
    \node[block, below=of ghat] (mult) {multiplicación por $\tau(\sigma)$\\(discreto: $M_\tau$)};
    \node[data, left=of mult] (gfhat) {$\widehat g_{\mathrm{filt}}$};
    \node[block, left=of gfhat] (ifft) {Fourier inversa\\$\mathcal F_t^{-1}$};
    \node[data, left=of ifft] (gf) {$g_{\mathrm{filt}}$};
    \node[block, below=of gf] (bp) {retroproyección\\$\mathcal R^*$};
    \node[data, right=of bp] (rec) {reconstrucción\\$f_\tau$};

    \draw[->, thick] (sino) -- (fft);
    \draw[->, thick] (fft) -- (ghat);
    \draw[->, thick] (ghat) -- (mult);
    \draw[->, thick] (mult) -- (gfhat);
    \draw[->, thick] (gfhat) -- (ifft);
    \draw[->, thick] (ifft) -- (gf);
    \draw[->, thick] (gf) -- (bp);
    \draw[->, thick] (bp) -- (rec);

    \node[align=center, below=0.45cm of mult]
    {rampa canónica continua: $\tau(\sigma)=|\sigma|$};
\end{tikzpicture}
\shorthandon{<>}
}
\caption{Flujo de la retroproyección filtrada en geometría paralela. Elaboración propia.}
\label{fig:fbp-flujo}
\end{figure}

La retroproyección filtrada es el algoritmo analítico clásico para datos de proyección paralela. Una presentación computacional se encuentra en Natterer \cite[Sec.~V.1, p.~102 y ss.]{natterer2001}; Kak y Slaney discuten la reconstrucción desde proyecciones paralelas en \cite[Sec.~3.3]{kak_slaney1988}. En esta tesis, la rampa se mantendrá como referencia teórica y computacional, mientras que las comparaciones numéricas posteriores reemplazarán $|\sigma|$ por otros símbolos $\tau(\sigma)$.
